\section{第一章 绪论}

\section{研究背景及意义}

（一）

随着我国经济的快速发展和城市化进程的不断推进，城市轨道交通已经成为现代城市的重要基础设施和公共交通体系的核心组成部分。城市轨道交通作为缓解城市交通拥堵、改善居民出行条件的关键解决方案，不仅承载着数亿乘客的日常出行需求，更承载着城市经济社会发展的重要职能。根据国家发改委公布的数据显示，截至2023年底，我国已有50个城市开通运营城市轨道交通线路，总运营里程超过9500公里，日均客运量超过5000万人次。这一庞大的交通网络规模不仅体现了我国城市轨道交通事业的蓬勃发展，也对供电系统的可靠性、安全性和稳定性提出了前所未有的高要求。

在城市轨道交通供电系统中，变电所扮演着至关重要的角色。它作为整个供电系统的核心枢纽，承担着将高压电能转换为适用于轨道交通牵引供电的电压等级的重要任务。变电所的稳定运行直接决定了轨道交通列车的供电质量和运行效率，而高压开关柜作为变电所的关键电气设备，其性能优劣更是影响着整个供电系统的可靠性。高压开关柜在变电所中主要承担着控制和保护电力系统的正常运行、实现电气设备的可靠接通和断开、提供故障隔离和系统保护功能等多项重要职责。任何高压开关柜的故障都可能导致供电中断，进而影响列车的正常运行，造成严重的经济损失和社会影响。

\begin{figure}[h]
\centering
%% % \includegraphics[width=0.8\textwidth]{figs/railway_power_system.png}
\caption{城市轨道交通供电系统架构示意图}
\label{fig:railway_power_system}
\end{figure}

（一）

近年来，随着城市轨道交通网络的不断扩大和设备运行年限的逐步增加，高压开关柜故障事件呈现出明显的上升趋势。根据相关统计数据显示，2022年全国城市轨道交通供电系统发生的各类故障中，高压开关柜相关故障占比超过25\%，已经成为制约轨道交通供电系统安全稳定运行的主要问题之一。这些故障不仅频率较高，而且往往具有突发性和破坏性强的特点。

高压开关柜故障主要可以分为几种类型：绝缘故障、机械故障、温升故障以及其他辅助系统故障。绝缘故障主要表现为绝缘击穿、局部放电等现象，通常是由于绝缘材料老化、环境湿度过高或电气应力过大等原因引起的；机械故障主要涉及操动机构卡涩、分合闸失败等问题，往往与机械部件磨损、润滑不足或制造质量有关；温升故障则表现为触头过热、连接部位异常发热等现象，多由接触电阻增大或负荷过大导致；其他故障则包括辅助回路故障、控制回路异常等问题。这些故障一旦发生，不仅会直接影响列车正常运行，还可能引发连锁反应，严重时会导致大面积停运，造成巨大的经济损失和社会影响。

\begin{figure}[h]
\centering
%% % \includegraphics[width=0.8\textwidth]{figs/switchgear_fault_types.png}
\caption{高压开关柜常见故障类型分布}
\label{fig:switchgear_fault_types}
\end{figure}

（一）

传统的离线检测方式虽然在一定程度上能够发现设备故障隐患，但存在诸多明显的局限性。首先，检测周期相对较长，无法实现对设备运行状态的实时掌握，往往需要在设备停运期间进行检测，这限制了检测的及时性；其次，受天窗时间限制，难以对所有设备进行全面细致的检测，很多潜在问题无法被及时发现；此外，传统的检测方式需要大量的人力和物力投入，检测费用高昂，经济性较差；最重要的是，这种检测方式无法对故障发展趋势进行预测，往往只能在故障发生后进行被动处理。

在线监测技术作为一种先进的设备状态监测手段，能够实现对高压开关柜运行状态的连续监测和实时分析，具有传统检测方式无法比拟的优势。在实时性方面，在线监测系统能够实现24小时不间断监测，确保设备状态随时可知；在准确性方面，通过多参数综合分析，可以显著提高故障诊断的精度和可靠性；在经济性方面，虽然初期投资相对较高，但长期来看能够大幅降低维护成本，提高设备利用率；在预测性方面，基于大数据分析和人工智能技术，可以实现故障预警，提前采取预防措施，避免故障发生。

\begin{figure}[h]
\centering
%% % \includegraphics[width=0.8\textwidth]{figs/monitoring_advantage.png}
\caption{在线监测技术优势对比分析}
\label{fig:monitoring_advantage}
\end{figure}

（一）

本研究以城市轨道交通变电所高压开关柜为研究对象，系统性地探讨在线监测技术的应用与维护策略优化，具有重要的理论意义和实践价值。

从理论意义上看，本研究丰富了城市轨道交通供电系统在线监测的理论体系，为高压电气设备状态监测技术在轨道交通领域的应用提供了重要的理论参考和实践指南。同时，本研究将人工智能、大数据分析等先进技术引入到电力设备维护领域，推动了传统电气设备维护模式的创新和发展，为相关学科的理论发展做出了积极贡献。

从实践意义上看，本研究能够显著提高城市轨道交通供电系统的运行可靠性，通过在线监测技术的应用，可以提前发现设备故障隐患，降低故障发生率，保障供电系统的稳定运行。同时，通过优化维护策略，可以有效降低设备维护成本和故障损失，为轨道交通运营企业带来显著的经济效益。此外，本研究为轨道交通行业制定维护标准和规范提供了重要的技术支撑和决策依据，有助于提升整个行业的维护水平。最后，本研究的应用成果将有力保障城市轨道交通的安全运营，为广大乘客的出行安全提供更加坚实的保障。

\section{国内外研究现状}

（一）

在线监测技术起源于20世纪80年代的电力系统状态监测研究。早期主要应用于大型发电机组的振动监测和温度监测。随着传感器技术和计算机技术的发展，在线监测技术逐渐扩展到高压电气设备领域。

国外在线监测技术起步较早，主要代表性研究包括：
\begin{enumerate}
    \item 美国西屋电气公司的机械特性在线监测系统
    \item 德国西门子公司的局放在线监测技术
    \item 日本三菱公司的温度场分布监测系统
\end{enumerate}

国内在线监测技术研究始于20世纪90年代，经过三十多年的发展，已形成较为完整的理论体系和技术体系。在城市轨道交通领域的应用研究主要集中在：
\begin{enumerate}
    \item 北京交通大学的牵引供电系统在线监测技术
    \item 上海交通大学的城市轨道交通设备状态监测平台
    \item 重庆大学的智能电网在线监测技术
\end{enumerate}

（一）

1.

温度监测是高压开关柜在线监测的基础技术，主要监测触头温度、母线温度等关键部位。当前主流技术包括：
\begin{enumerate}
    \item 红外测温技术：非接触式测量，适用于运行中设备
    \item 光纤测温技术：抗电磁干扰能力强，适用于强电场环境
    \item 热敏电阻测温技术：成本低廉，易于安装
\end{enumerate}

1.

局部放电是绝缘劣化的重要标志，局放监测技术包括：
\begin{enumerate}
    \item 超声波法：检测放电产生的超声波信号
    \item 高频法：检测高频电流脉冲信号
    \item 特高频法：检测特高频电磁波信号
\end{enumerate}

1.

机械特性监测主要评估开关柜的分合闸性能，包括：
\begin{enumerate}
    \item 行程特性监测：监测动触头行程曲线
    \item 速度特性监测：监测分合闸速度变化
    \item 力特性监测：监测操动力和触头压力
\end{enumerate}

（一）

1.

传统的高压开关柜维护主要采用定期检修模式，即按照固定周期进行预防性维护。这种模式存在以下问题：
\begin{enumerate}
    \item 维护过度：对正常设备进行不必要的检修
    \item 维护不足：无法及时发现潜在故障
    \item 资源浪费：人力、物力、财力浪费严重
\end{enumerate}

1.

状态检修（CBM）基于设备实际运行状态进行维护决策，主要包括：
\begin{enumerate}
    \item 故障预测：基于历史数据预测故障发生时间
    \item 风险评估：评估设备故障对系统的影响程度
    \item 维护优化：确定最佳维护时机和维护内容
\end{enumerate}

1.

预测性维护结合人工智能技术，实现智能维护决策：
\begin{enumerate}
    \item 机器学习算法：故障模式识别和预测
    \item 大数据分析：多源数据融合分析
    \item 专家系统：基于规则的智能诊断
\end{enumerate}

（一）

1.

美国纽约地铁采用的在线监测系统实现了对500多台高压开关柜的实时监测，故障预测准确率达到95\%，每年节省维护成本约200万美元。

德国柏林地铁的智能维护系统将传统定期检修周期从6个月延长到18个月，设备可用性提高了15\%。

1.

国内城市轨道交通在线监测技术应用尚处于起步阶段，主要集中在：
\begin{enumerate}
    \item 北京地铁：试点应用局放和温度在线监测
    \item 上海地铁：建设综合状态监测平台
    \item 深圳地铁：开展机械特性在线监测研究
\end{enumerate}

但总体来看，国内在线监测技术应用深度不够，系统集成度不高，维护策略智能化水平有待提高。

（一）

尽管在线监测技术取得了长足发展，但在城市轨道交通变电所高压开关柜领域的应用仍存在诸多问题：

1.
\begin{enumerate}
    \item 传感器技术不成熟：抗干扰能力弱，长期稳定性差
    \item 数据处理能力不足：海量监测数据难以有效分析
    \item 系统集成度不高：各监测子系统相对独立
\end{enumerate}

1.
\begin{enumerate}
    \item 缺乏统一的技术标准和规范
    \item 维护人员技术水平参差不齐
    \item 设备更新改造成本高昂
\end{enumerate}

1.
\begin{enumerate}
    \item 缺乏有效的维护策略体系
    \item 数据共享机制不完善
    \item 决策支持系统不够完善
\end{enumerate}

\section{本文研究目标和研究内容}

（一）

本研究旨在针对城市轨道交通变电所高压开关柜的特点和运行环境，构建完整的在线监测技术体系和智能维护策略体系，实现以下目标：
\begin{enumerate}
    \item 建立高压开关柜在线监测技术框架
    \item 开发多参数综合监测系统
    \item 提出基于风险评估的预防性维护策略
    \item 构建智能故障诊断和预测模型
    \item 制定维护策略优化方案
\end{enumerate}

（一）

本文主要研究内容包括以下三个方面：

\begin{enumerate}
    \item \textbf{高压开关柜在线监测技术研究}：分析高压开关柜故障机理，研究温度监测、局放监测、机械特性监测等关键技术，建立多参数综合监测体系，探讨大数据分析在故障预警中的应用。

    \item \textbf{维护策略分析与优化}：基于在线监测数据，建立设备状态评估模型，提出状态检修与定期检修相结合的混合维护模式，构建基于风险评估的预防性维护体系。

    \item \textbf{系统集成与应用验证}：设计在线监测系统架构，开发监测软件平台，通过实际案例验证系统效能，提出技术推广建议和标准化方案。
\end{enumerate}